\begin{abstract}
{ 
Value nodes with elementary operations can be built to create complex
functions that achieve practical appliances. Although more memory intensive,
this serves as a bare minimum for artificial intelligence architectures---that is 
without the computational optimizations. These value nodes hold a history of past 
operations and children nodes, creating a directed acyclic graph that allows us to 
calculate local gradients. Through these value nodes, we build up neurons, layers, and
all the way up to a multilayer perceptron. With the MNIST dataset, we're able demonstrate the functionality
of the value node multilayer perceptron using a training loop with gradient descent
resulting in empirical results of 92$\sim$\% accuracy. Utilizing optimizations such as L2 normalization,
adaptive moment estimation (Adam), cross entropy loss, and Gaussian error linear unit (GELU) resulted
to lower loss and a final accuracy of 97$\sim$\%.
\vspace{4ex}}
\end{abstract}




